\documentclass[11pt,a4paper]{article}

\usepackage[T1]{fontenc}
\usepackage[utf8]{inputenc}
\usepackage[brazil]{babel}
\usepackage{lmodern}
\usepackage{microtype}
\usepackage[margin=2cm]{geometry}
\usepackage{hyperref}
\usepackage{titlesec}
\usepackage{xcolor}

\hypersetup{
  colorlinks=true,
  linkcolor=black,
  urlcolor=black,
  citecolor=black,
  pdfborder={0 0 0}
}

\definecolor{sectioncolor}{HTML}{111111}
\definecolor{accentcolor}{HTML}{1A1A1A}

\pagestyle{empty}
\setlength{\parindent}{0pt}
\setlength{\parskip}{5pt}
\linespread{1.05}
\selectfont

\titleformat{\section}
  {\bfseries\large\color{sectioncolor}}
  {}{0pt}{}[\vspace{-2pt}{\color{sectioncolor}\titlerule[0.8pt]}]
\titlespacing*{\section}{0pt}{12pt}{8pt}

\newcommand{\cvrole}[4]{%
  \vspace{4pt}%
  {\bfseries\color{accentcolor}#1}\hfill{\itshape #2}\\[-2pt]%
  {\itshape #3}\hfill #4\\[6pt]%
}

\begin{document}

\begin{center}
  {\fontsize{20}{24}\selectfont\bfseries Kamille Hartmann Maccagnan}\\[8pt]
  {\large Analista de Processos \textbar{} Eficiência Operacional \textbar{} Power BI}\\[10pt]
  Campo Comprido, Curitiba-PR \enspace\textbullet\enspace (41) 9 9928-4424 \enspace\textbullet\enspace
  \href{mailto:kamillehartmannm@gmail.com}{kamillehartmannm@gmail.com}\\[3pt]
  \href{https://linkedin.com/in/kamille-hartmann-maccagnan/}{linkedin.com/in/kamille-hartmann-maccagnan/}
\end{center}

\vspace{6pt}

\section{Resumo Profissional}

Profissional com experiência em mapeamento e padronização de processos, automação operacional e monitoramento de indicadores em ambientes de tecnologia, varejo e operações industriais. Atua na documentação de processos, elaboração de POPs, implementação de automações com Power Automate e desenvolvimento de dashboards em Power BI para acompanhamento de KPIs e apoio à tomada de decisão. Vivência na automação de distribuição de relatórios de BI e no uso de recursos de linguagem natural do Power BI para consulta a indicadores. Experiência em projetos de melhoria contínua, gestão de demandas e engajamento de stakeholders entre áreas de negócio e tecnologia. Perfil analítico, organizado e orientado a dados, com capacidade de estruturar problemas e conduzir iniciativas de eficiência operacional.

\section{Experiência Profissional}

\cvrole{Anjuss}{Analista de Suporte de TI}{07/2025 -- Atual}

Atuo no mapeamento, padronização e melhoria de processos operacionais, com criação e atualização de documentações, POPs para a rede de lojas e POPs voltados à equipe de T.I., assegurando qualidade, consistência e confiabilidade das rotinas executadas. Lidero treinamentos e capacitação de novos funcionários, promovendo adoção de ferramentas e práticas padronizadas entre as áreas. Desenvolvo dashboards em Power BI para estruturação e acompanhamento de indicadores de desempenho, apoiando a identificação de oportunidades de melhoria e priorização de ações. Implemento automações de processos com Power Automate, incluindo gatilhos, fluxos e integrações entre rotinas analíticas e sistemas, aumentando eficiência operacional e reduzindo retrabalho.

\cvrole{HORSE do Brasil}{Estagiária de TI (IS/IT LATAM)}{07/2024 -- 07/2025}

Atuei em projetos de melhoria contínua, PMO e governança de dados em ambiente multinacional industrial, com acompanhamento de portfólio de projetos, cronogramas e entregas. Participei do mapeamento e documentação de processos, análise de riscos e padronização de informações, apoiando a evolução de práticas de governança. Elaborei relatórios executivos e dashboards em Power BI para monitoramento de KPIs de projetos e iniciativas estratégicas, com apoio ao uso do recurso de Perguntas e Respostas para consulta a indicadores por meio de linguagem natural. Implementei automações com Power Automate para distribuição recorrente de relatórios de BI, incluindo gatilhos, fluxos e integrações entre rotinas analíticas e sistemas. Utilizei o ServiceNow para gestão de demandas e atuei como ponte entre áreas de TI, operações, fornecedores e stakeholders internacionais, promovendo alinhamento e adoção de melhorias em contextos com múltiplas prioridades.

\cvrole{Restaurante Familiar}{Analista de Dados e Suporte Técnico}{01/2023 -- 07/2024}

Atuei no diagnóstico de processos operacionais e na análise de performance, vendas e resultados financeiros para apoio à gestão do negócio. Estruturei indicadores de desempenho e dashboards em Power BI para acompanhamento de resultados, identificando desvios e oportunidades de melhoria com impacto na eficiência operacional. Realizei padronização e validação de dados para garantir confiabilidade das informações utilizadas em relatórios e decisões. Atuei no alinhamento entre áreas operacionais e administrativas, contribuindo para simplificação de rotinas e melhoria da operação do dia a dia.

\section{Competências Técnicas}

\textbf{Processos e Eficiência Operacional:} Mapeamento de processos, padronização de processos, documentação operacional, POPs, melhoria contínua, gestão de demandas, diagnóstico operacional, Boas Práticas PMOBOK, Scrum, Kanban.

\textbf{Indicadores e Dados:} KPIs, monitoramento de desempenho, dashboards em Power BI (DAX, Power Query), Excel avançado, SQL, análise de dados, validação e qualidade de dados, relatórios gerenciais.

\textbf{Automação e Tecnologia:} Automação de processos (Power Automate), gatilhos e fluxos de trabalho, distribuição automatizada de relatórios de BI, integração entre sistemas e rotinas analíticas, recursos de linguagem natural do Power BI (Perguntas e Respostas), ServiceNow, Power BI.

\textbf{Gestão e Stakeholders:} Engajamento de stakeholders, comunicação com lideranças, interface entre negócio e tecnologia, treinamento e capacitação, gestão de múltiplas demandas, projetos multidisciplinares.

\textbf{Ferramentas:} Power BI, Power Automate, ServiceNow, Trello, Monday, Google Sheets, Pacote Office.

\textbf{Idiomas:} Inglês Avançado \enspace\textbullet\enspace Espanhol Básico \enspace\textbullet\enspace Coreano Básico

\section{Projetos e Conquistas}

1º lugar no Transformation Day Renault 2023 com solução orientada a melhoria de processos e dados. Participação no Hackathon Bosch 2023 com foco em resolução de problemas de negócio. Automação da distribuição de relatórios de BI com Power Automate e apoio ao uso de Perguntas e Respostas no Power BI na HORSE. Implementação de automações com Power Automate integradas ao Power BI para padronização de fluxos e redução de retrabalho. Estruturação de KPIs e dashboards para monitoramento de performance operacional. Desenvolvimento de documentações e POPs para padronização de rotinas em rede de lojas.

\section{Formação Acadêmica}

\textbf{PUCPR} \hfill Curitiba, PR\\
Graduação em Gestão da Tecnologia da Informação \hfill Previsão de conclusão: 06/2026\\[6pt]

\textbf{Escola Conquer} \hfill Curitiba, PR\\
Pós-graduação em Liderança e Gestão em Tecnologia \hfill Conclusão: 2024\\[6pt]

\textbf{Universidade Tuiuti do Paraná} \hfill Curitiba, PR\\
Graduação em Medicina Veterinária \hfill 06/2019\\[6pt]

\textbf{Qualittas} \hfill Curitiba, PR\\
Pós-graduação em Clínica Médica e Cirúrgica de Animais Exóticos e Pets Não Convencionais \hfill 12/2022

\end{document}
